\section*{Chapter \ref{ch:b2:fluid-kinematics} --- Fluid Kinematics}

\begin{solution}{exo:b2:fluid-kinematics:1}
(a) $\ell/L \sim 7 \times 10^{-8}$: continuum. (b) $0.07$: marginal --- the
continuum laws start to fail (slip at the walls). (c) $\ell = 70 \times 10^{-9}
/3 \times 10^{-7} = \qty{0.23}{m}$, $\ell/L = 0.23$: no continuum; the gas is
a rain of independent molecules (free-molecular regime).
\end{solution}

\begin{solution}{exo:b2:fluid-kinematics:2}
$a = v\,\dd v/\dd x = v_0^2(1 + x/L)/L$ (from $v_0^2/L$ to $2v_0^2/L$).
$\dd x/\dd t = v_0(1 + x/L)$: $t = (L/v_0)\ln2 = 0.69\,L/v_0$, between
$L/2v_0$ and $L/v_0$.
\end{solution}

\begin{solution}{exo:b2:fluid-kinematics:3}
$D_V = \int_0^Rv_{\max}(1 - r^2/R^2)\,2\pi r\,\dd r = \tfrac12\pi R^2v_{\max}$;
$\langle v\rangle = v_{\max}/2 = \qty{0.20}{m/s}$; $D_V = \qty{1.6e-5}{m^3/s} =
\qty{16}{mL/s}$.
\end{solution}

\begin{solution}{exo:b2:fluid-kinematics:4}
(a) $\operatorname{div} = 0$, $\vect\omega = \vect 0$: incompressible,
irrotational. (b) $\operatorname{div} = 0$, $\omega_z = 2\Omega$: solid
rotation. (c) $\operatorname{div} = 2a \ne 0$, irrotational (a plane
source). (d) $\operatorname{div} = 0$, $\omega_z = -k$: shear. (e)
$\operatorname{div} = b$: incompressible iff $b = 0$; irrotational for
every $b$.
\end{solution}

\begin{solution}{exo:b2:fluid-kinematics:5}
$D_V = \qty{0.20}{L/s}$. Hose $S = \qty{1.13}{cm^2}$: $v = \qty{1.8}{m/s}$; nozzle
$S = \qty{7.1}{mm^2}$: $v = \qty{28}{m/s}$. Mean acceleration $\Delta(v^2/2)/L
= (28^2 - 1.8^2)/0.08 \approx \qty{1e4}{m/s^2} \approx 1000\,g$. Time $\approx
\int\dd x/v$: a few milliseconds (\qty{4}{ms} for a linear speed profile).
\end{solution}

\begin{solution}{exo:b2:fluid-kinematics:6}
(a) $\dd y/\dd x = (V/U)\cos\omega t$: straight lines of slope
$(V/U)\cos\omega t$, all parallel, tilting with time. (b) $x = Ut$, $y =
(V/\omega)\sin\omega t$: the sinusoid $y = (V/\omega)\sin(\omega x/U)$. (c) At
$t = 0$ the \omterm{def:b2:fluid-kinematics:euler}{streamlines} are at \qty{45}{\degree} and the \omterm{def:b2:fluid-kinematics:euler}{pathline} leaves
the origin at \qty{45}{\degree}; at $t = \pi/2\omega$ the \omterm{def:b2:fluid-kinematics:euler}{streamlines} are
horizontal while the \omterm{def:b2:fluid-kinematics:euler}{pathline} is a fixed sinusoid. They differ because
the flow is unsteady: a \omterm{def:b2:fluid-kinematics:euler}{streamline} is a snapshot, a \omterm{def:b2:fluid-kinematics:euler}{pathline} a history.
(d) $\vect a = \partial_t\vect v = (0, -V\omega\sin\omega t)$ --- the field is
uniform, so the convective term vanishes.
\end{solution}

\begin{solution}{exo:b2:fluid-kinematics:7}
(a) $D_V = 50 \times 2 \times 1.2 = \qty{120}{m^3/s}$. (b) $120/60 = \qty{2.0}{m/s}$.
(c) $150/150 = \qty{1.0}{m/s}$. (d) $\rho \propto 1/T$ at constant $P$ halves;
$\rho vS$ is conserved: the speed doubles.
\end{solution}

\begin{solution}{exo:b2:fluid-kinematics:8}
(a) $\omega_z = \frac1r\frac{\dd(rv_\theta)}{\dd r}$: $2\Omega$ inside, $0$
outside. (b) Inside $\Gamma = 2\pi r\cdot\Omega r = 2\Omega\cdot\pi r^2$, the flux
of the \omterm{def:b2:fluid-kinematics:vorticity}{vorticity}; outside $\Gamma = 2\pi\Omega a^2$, constant, the flux of
the whole core. (c) $\Omega = v_{\max}/a = \qty{1.4}{rad/s}$; $\Gamma = 2\pi av_{\max}
= \qty{2.2e4}{m^2/s}$; at \qty{500}{m}: $\Omega a^2/r = \qty{7}{m/s}$. (d) The
core rotates like a solid; outside, the flow is irrotational: particles
go round the axis without spinning about themselves. The \omterm{def:b2:fluid-kinematics:vorticity}{vorticity} ---
the rotation --- lives in the core.
\end{solution}

\begin{solution}{exo:b2:fluid-kinematics:9}
(a) $\varphi = \tfrac12k(x^2 - y^2)$, $\Delta\varphi = k - k = 0$. (b) $\dd x/
kx = \dd y/(-ky)$: $xy = $ const; $\partial_y\psi = kx = v_x$, $-\partial_x\psi =
-ky = v_y$. (c) $\vect a = (v_x\partial_x + v_y\partial_y)\vect v = k^2(x, y) =
\operatorname{\vect{grad}}\bigl(\tfrac12k^2(x^2 + y^2)\bigr) = \operatorname{\vect{grad}}
(v^2/2)$: in a steady \omterm{def:b2:fluid-kinematics:vorticity}{irrotational flow} the acceleration is the
gradient of the kinetic energy per unit mass; it points away from the
stagnation point --- a particle decelerates as it approaches the wall
along $y$ and accelerates away along $x$. (d) $\dot x = kx$, $\dot y = -ky$:
$x = x_0\eu^{kt}$, $y = y_0\eu^{-kt}$ ($xy$ constant).
\end{solution}

\begin{solution}{exo:b2:fluid-kinematics:10}
(a) $\dd x/\dd t = x/(t + \tau)$ gives $x = x_0(1 + t/\tau)$, $\ddot x = 0$;
\omterm{thm:b2:fluid-kinematics:material}{material derivative}: $\partial_tv + v\partial_xv = -x/(t + \tau)^2 + x/(t + \tau)^2
= 0$. (b) $\partial_t\rho + \partial_x(\rho v) = \dot\rho + \rho/(t + \tau) = 0$: $\rho =
\rho_0\tau/(t + \tau)$. (c) Two particles' separation grows as $(1 + t/\tau)$
and the mass between them is fixed: $\rho \propto 1/(1 + t/\tau)$. (d) $v =
Hx$ with $H = 1/(t + \tau)$: a Hubble law, free expansion (no
acceleration), density falling as the inverse of the scale factor.
\end{solution}

\begin{solution}{exo:b2:fluid-kinematics:11}
(a) $\vect v = \operatorname{\vect{grad}}\varphi = (U + qx/2\pi r^2, qy/2\pi r^2)$;
the source has $v_r = q/2\pi r$, $\operatorname{div} = \frac1r\dd(rv_r)/\dd r
= 0$ and no $\theta$ dependence, so $\omega_z = 0$; sums of such fields are
again incompressible and irrotational. (b) On the negative $x$ axis,
$U - q/2\pi r = 0$: $x_s = -q/2\pi U$. (c) At the stagnation point $\theta =
\pi$, $y = 0$: $\psi = q/2$; the \omterm{def:b2:fluid-kinematics:euler}{streamline} $Ur\sin\theta + q\theta/2\pi = q/2$,
i.e.\ $r\sin\theta = (q/2\pi U)(\pi - \theta)$ --- a curve leaving $x_s$ and
opening downstream. (d) $\theta \to 0$: $y \to q/2U$; $\theta \to 2\pi$ (lower
branch): $y \to -q/2U$. The fluid inside that \omterm{def:b2:fluid-kinematics:euler}{streamline} comes from the
source; outside it, the uniform stream flows around a half-body of
asymptotic width $q/U$.
\end{solution}

\begin{solution}{exo:b2:fluid-kinematics:12}
(a) $\operatorname{div}\vect v = \frac1r\partial_r(-q/2\pi) = 0$; $\omega_z =
\frac1r\partial_r(\Gamma/2\pi) = 0$. (b) $\dd r/v_r = r\,\dd\theta/v_\theta$ gives
$\dd r/r = -(q/\Gamma)\dd\theta$. (c) $a_r = v_r\partial_rv_r - v_\theta^2/r =
-(q^2 + \Gamma^2)/4\pi^2r^3 = -v^2/r$; $a_\theta = v_r\partial_rv_\theta + v_rv_\theta/r
= 0$. (d) $\dd r/\dd t = -q/2\pi r$: $t = \pi(r_0^2 - a^2)/q$; turns
$= (\Gamma/2\pi q)\ln(r_0/a)$. (e) $t = \qty{520}{s} \approx \qty{9}{min}$; $10$
turns.
\end{solution}

\begin{solution}{pb:b2:fluid-kinematics:1}
\textbf{1.} Through a cylinder of radius $r$: $2\pi rh|v_r| = D_V$, so
$v_r = -q/2\pi r$, $q = D_V/h = \qty{1.5e-3}{m^2/s}$.

\textbf{2.} $rv_\theta = 0.5 \times 0.01 = \qty{5e-3}{m^2/s}$: $v_\theta = \Gamma/2\pi r$
with $\Gamma = \qty{3.1e-2}{m^2/s}$.

\textbf{3.} $\operatorname{div}\vect v = \frac1r\partial_r(rv_r) = 0$, $\omega_z =
\frac1r\partial_r(rv_\theta) = 0$.

\textbf{4.} $\dd r/r = -(q/\Gamma)\dd\theta$: $r = r_0\eu^{-q\theta/\Gamma}$; turns
$= (\Gamma/2\pi q)\ln(r_0/a) = 3.3 \times 3.2 \approx 11$.

\textbf{5.} $t = \pi(r_0^2 - a^2)/q = \qty{520}{s} \approx \qty{9}{min}$; $v_\theta(a) =
\qty{0.25}{m/s}$.

\textbf{6.} $\vect a = -(v^2/r)\vect e_r$; at $r = a$: $v^2 = 0.25^2 + 0.012^2$,
$a = \qty{3.1}{m/s^2} \approx 0.3g$.

\textbf{7.} $\Omega_\oplus\sin\lambda\,r = 5 \times 10^{-5} \times 0.5 = \qty{2.5e-5}{m/s}$,
four hundred times less than the stirring: the hemisphere decides
nothing in an ordinary bath (it does in a tank left still for a day).

\textbf{8.} $v_\theta = \Omega r$ ($r < a$), $\Omega a^2/r$ ($r > a$); $\Omega = 80/60
= \qty{1.33}{rad/s}$, $\Gamma = 2\pi av_{\max} = \qty{3.0e4}{m^2/s}$.

\textbf{9.} $\omega_z = 2\Omega = \qty{2.7}{s^{-1}}$ in the core, $0$ outside;
$v_\theta$ rises linearly to \qty{80}{m/s} then falls as $1/r$; $\omega_z$ is a
step.

\textbf{10.} $a = v_{\max}^2/a = 6400/60 = \qty{107}{m/s^2} \approx 11g$, inward.

\textbf{11.} Outside, the \omterm{def:b2:fluid-kinematics:vorticity}{vorticity} is zero: a small element translates
along its circle without rotating about itself --- the paddle wheel
keeps its orientation. In the core (solid rotation) it turns with the
fluid, once per revolution.

\textbf{12.} $E_{\text{core}} = \int_0^a\tfrac12\rho\Omega^2r^2\,2\pi r\,\dd r =
\tfrac14\pi\rho\Omega^2a^4 = \qty{2.2e7}{J/m}$.

\textbf{13.} $E_{\text{out}} = \int_a^R\tfrac12\rho(\Omega a^2/r)^2\,2\pi r\,\dd r =
\pi\rho\Omega^2a^4\ln(R/a) = 8.7 \times 10^7 \times 4.4 = \qty{3.8e8}{J/m}$; the
integrand $\rho v^2r \propto 1/r$, hence the logarithm.

\textbf{14.} $\qty{4.0e8}{J/m} \times \qty{1000}{m} = \qty{4e11}{J}$: about
\qty{100}{t} of TNT.

\textbf{15.} Inflow $2\pi a \times 300 \times 5 = \qty{5.7e5}{m^3/s}$ through
$\pi a^2 = \qty{1.1e4}{m^2}$: $v_z = \qty{50}{m/s}$.

\textbf{16.} $\omega = 2v_{\max}/a = \qty{2.5e-3}{s^{-1}} = 50f$; $\Gamma = 2\pi av_{\max}
= \qty{1.3e7}{m^2/s}$.

\textbf{17.} $r_0\cdot fr_0/2 = r(v_\theta + fr/2)$: $v_\theta = f(r_0^2 - r^2)/2r$;
\qty{60}{m/s} at \qty{100}{km}, \qty{155}{m/s} at \qty{40}{km}.

\textbf{18.} $f > 0$ in the northern hemisphere gives $v_\theta > 0$:
counterclockwise seen from above (cyclonic). Over \qty{500}{km} the
planetary rotation is the dominant initial rotation; in a bath the
stirring is.

\textbf{19.} Friction with the sea removes angular momentum, and the
pressure gradient cannot balance the centripetal acceleration of an
ever faster ring: the air rises before reaching the axis, in the eye
wall, leaving a calm eye.

\textbf{20.} $\rho\,2\pi rHv = 1.2 \times 2\pi \times 5 \times 10^5 \times 10^3 \times 5 =
\qty{1.9e10}{kg/s}$.

\textbf{21.} $\vect v(M, t) = -U\vect e_x + \vect v_0(M + Ut\,\vect e_x)$: the
pattern is advected, the field is unsteady in the ground frame.
The translation adds to the swirl on the side where $v_\theta$ points
west --- the north (right-hand) side of the westward track.

\textbf{22.} $0.05 \times 10^7/3600 = \qty{139}{m^3/s}$; $v = 139/40 = \qty{3.5}{m/s}$.

\textbf{23.} $D_V = \qty{83}{cm^3/s}$: aorta \qty{0.33}{m/s}; capillaries
\qty{0.33}{mm/s}; \qty{3}{s} per millimetre.

\textbf{24.} $S\,\dd h/\dd t = -s\sqrt{2gh}$: $\sqrt h = \sqrt{h_0} -
\frac sS\sqrt{g/2}\,t$, empty at $T = (S/s)\sqrt{2h_0/g} = 10^4 \times 0.45 =
\qty{4500}{s} \approx \qty{75}{min}$.

\textbf{25.} Bath: $r \sim \qty{0.5}{m}$, $v \sim \qty{1}{cm/s}$, $\Gamma =
\qty{3e-2}{m^2/s}$; tornado: \qty{60}{m}, \qty{80}{m/s}, $\omega = \qty{2.7}{s^{-1}}$,
$\Gamma = \qty{3e4}{m^2/s}$; cyclone: \qty{40}{km}, \qty{50}{m/s}, $\omega =
\qty{2.5e-3}{s^{-1}}$, $\Gamma = \qty{1.3e7}{m^2/s}$. In each, the angular
momentum per unit mass $rv_\theta$ of the converging fluid is conserved:
inflow spins it up.
\end{solution}
